\documentclass[10pt,a4paper]{amsart}
\usepackage[margin=19mm]{geometry}
\usepackage{amsmath,amssymb,mathtools}
\usepackage[scr=rsfs,cal=euler]{mathalfa}
\usepackage{xfrac}
\usepackage[unicode,hidelinks,bookmarks=false]{hyperref}
\newcommand{\ie}{i.e.\ }
\newcommand{\cf}{cf.\ }
\newcommand{\resp}{respectively\ }
\newcommand{\Subsection}{\subsection}
\providecommand{\nobreakdash}{\nobreak\hbox{-}}
\setlength{\emergencystretch}{2em}
\multlinegap0pt
\usepackage{mathtools}
\usepackage[shortlabels]{enumitem}
\usepackage{xcolor}
\usepackage{amssymb}
\usepackage{tikz}
\newtheorem{lemma}{Lemma}
\title{Skolem-Mahler-Lech in rings of positive characteristic: a shorter proof and a multi-dimensional generalization}
\author{Ruiwen Dong, Doron Shafrir}
\date{Source: arXiv:2609.03127v1 (CC BY 4.0)}
\begin{document}
\maketitle

\noindent\emph{Source and licence.} Skolem-Mahler-Lech in rings of positive characteristic: a shorter proof and a multi-dimensional generalization. Version arXiv:2609.03127v1 (CC BY 4.0), distributed by arXiv under the Creative Commons Attribution 4.0 International licence, http://creativecommons.org/licenses/by/4.0/, as recorded in the arXiv licence metadata for this exact version and shown on the abstract page https://arxiv.org/abs/2609.03127v1. Full licence notice: \texttt{licenses/arxiv-2609.03127-CC-BY-4.0.txt}.\par\smallskip

\noindent\textit{Editorial scope.} Counted unit: the article's lemma lem:split\_extension (source0.tex lines 1229--1231). The article's complete original proof of it is delivered with it (source0.tex lines 1232--1238, one \texttt{proof} environment of 7 lines). No mathematics is written, repaired or re-derived: the statement and the proof are the article's own lines, in the article's own notation, and no retained line is substituted or re-flowed.  No further source-local context is needed: every cross-reference of every retained line resolves inside the two delivered spans.  Nothing was omitted from any retained span, so the package carries no omission record.  No retained line contains a citation, so the wrapper carries no bibliography and no bibliography entry.  No retained line contains a figure environment, an image-inclusion command or drawing code, so no image is delivered and the package declares no asset dependency.  Editorial formatting only, in addition: an AMS \texttt{amsart} preamble that restates the article's own declarations and macros used by the retained lines, namely the environments \texttt{lemma} on separate counters, together with the standard packages those lines need.  The retained environments print in this document as Lemma 1 in order of appearance, whereas the article prints the same items with the numbering of its own full document.  The complete native source of the article as distributed in the arXiv e-print for this exact version is bundled unchanged as \texttt{sources/209266/source0.tex}.
\begin{lemma}[folklore]\label{lem:split_extension}
    Let $R$ be a ring, $M$ be an $R$-module, and $g \in R[T]$ be a monic polynomial. We can compute a ring extension $\widetilde R \supseteq R$, finitely generated as an $R$-module, such that $g$ splits over $\widetilde R$, and such that the map $M\rightarrow M\otimes _R \widetilde R$ is injective. 
\end{lemma}

\begin{proof}
    We prove by induction on the degree $d \coloneqq \deg(g)$. 
    Define $\widetilde R$ to be the quotient $R[x]/\langle g(x)\rangle$ of the polynomial ring $R[x]$. 
    As $R$-modules, $\widetilde R = R + xR + \cdots + x^{d-1}R \cong R^d$, and therefore $\widetilde M =M\otimes_R \widetilde R \cong M^d$, and the map $M\rightarrow M\otimes _R \widetilde R$ is injective. 
    In the ring $\widetilde R[T]$, we have $T - x \mid g(T)$.
    Therefore we can let $\widetilde g(T) \coloneqq \frac{g(T)}{T-x}$ and use induction on $\widetilde R, \widetilde M, \widetilde g$. 
\end{proof}

\end{document}
